\documentclass[a4paper]{article}

\usepackage{amsfonts, amsmath, amssymb, amsthm}
\usepackage{fullpage}
\usepackage{hyperref}

\newtheorem{theorem}{Theorem}[section]
\newtheorem{lemma}[theorem]{Lemma}
\newtheorem{corollary}[theorem]{Corollary}
\newtheorem{remark}[theorem]{Remark}

\newcommand{\F}{\mathbb{F}}
\newcommand{\abs}[1]{\left|#1\right|}
\DeclareMathOperator{\rank}{rank}
\renewcommand{\qedsymbol}{$\blacksquare$}

\setlength{\parindent}{0em}
\setlength{\parskip}{1em}

\begin{document}

\title{Finite Witness Bounds for \textit{Lights Out} Nullities}
\author{William Boyles}
\date{\today}
\maketitle

\begin{abstract}
	Let $d(n)$ be the nullity over $\F_2$ of the $n\times n$ bounded
	\textit{Lights Out} grid.
	We prove that for every positive integer $D$, if $D=d(n)$ for some
	$n$, then $D=d(n_0)$ for an explicitly bounded $n_0$.
	In particular, whether a given integer occurs as a grid nullity is
	decidable by a finite computation.
	Combining the witness bound with a published algorithm for evaluating
	$d(n)$ gives an exponential-time decision procedure in the numerical
	target nullity, or a double-exponential upper bound when the target is
	encoded in binary.
	The general bound is exponential and deliberately coarse, but the
	argument gives a much smaller bound for individual values.
	An exact computation of the resulting finite search proves that 34 is
	not a grid nullity.
\end{abstract}

\section{Motivation and prior work}

Prior work gives efficient methods for evaluating $d(n)$ at a specified
board size and devotes particular attention to reversibility, the condition
$d(n)=0$.
Goldwasser, Klostermeyer, and Trapp give an
$O(n\log^2 n)$ polynomial-GCD algorithm for the nullity of a specified
grid dimension \cite{GoldwasserKlostermeyerTrapp1997}.
Although their algorithm is framed as a reversibility test, it constructs
the full polynomial GCD, whose degree is the nullity, without changing the
asymptotic running time.
Sutner proves periodicity when one dimension of a rectangular grid is
fixed \cite{Sutner1988}, while Goldwasser, Klostermeyer, and Ware and
Hunziker, Machiavelo, and Park bound the Fibonacci indices of individual
irreducible factors \cite{Goldwasser2002,HUNZIKER2004465}.
Hunziker, Machiavelo, and Park also prove that there are infinitely many
primitive indices at which the square grid is irreversible.
These results concern evaluation, zero versus positive nullity, or
individual polynomial factors; they do not determine which positive
nullity values occur.

Define the image-membership problem
\[
	\operatorname{IMAGE}_d
	=
	\{N\geq0:\text{there exists }n\geq0\text{ with }d(n)=N\}.
\]
Since each $d(n)$ can be computed, this set is computably enumerable:
one may test $n=0,1,2,\ldots$ until a witness is found.
This procedure does not halt when $N$ is absent.
Thus pointwise computability of $d$ does not by itself imply that
$\operatorname{IMAGE}_d$ is decidable.
This distinction is familiar from linear recurrence sequences, where
term evaluation is effective but the general Skolem problem of deciding
whether a zero ever occurs remains open \cite{BiluEtAl2022}.
We do not claim a reduction from that problem; it merely illustrates why
an effective witness bound is substantive.

To our knowledge, prior work does not give a bound depending only on a
target nullity $N$, prove that $\operatorname{IMAGE}_d$ is decidable, or
give an explicit complexity upper bound for this image-membership problem.
The purpose of this paper is to supply such a bound.
Besides settling decidability, the bound turns finite computations into
certificates that particular values are globally impossible.
This is motivated by the first gaps in the observed range.
Every positive even value below 26 occurs \cite{OEISA159257}.
A companion manuscript at \path{tex/finite_fields/finite_fields.tex}
proves that 26 is impossible and thereby excludes the infinite family
$2^t(26+2)-2$.
In the computation through $n=100000$ recorded in
\path{code/serialization/b159257.txt}, the value 34 is the next gap not
belonging to that family.

\section{Preliminaries}

Let $F_n(x)\in\F_2[x]$ be the Fibonacci polynomials
\[
	F_0(x)=0,\qquad F_1(x)=1,\qquad
	F_n(x)=xF_{n-1}(x)+F_{n-2}(x).
\]
The nullity of the $n\times n$ grid is
\begin{equation}\label{eq:nullity}
	d(n)=
	\deg\gcd\!\left(F_{n+1}(x),F_{n+1}(x+1)\right)
\end{equation}
\cite{Sutner00sigma-automataand}.
We include the empty grid by setting $d(0)=0$, which also agrees with
\eqref{eq:nullity}.

We use the Fibonacci divisibility identity
\begin{equation}\label{eq:fibonacci-gcd}
	\gcd(F_m,F_n)=F_{\gcd(m,n)}
\end{equation}
over $\F_2[x]$ \cite{HUNZIKER2004465}.
For a non-zero polynomial $P\in\F_2[x]$, define
\[
	\rank(P)=\min\{r>0:P\mid F_r\}.
\]
Every non-zero polynomial has finite rank \cite{HUNZIKER2004465}.
Equation \eqref{eq:fibonacci-gcd} gives the rank criterion
\begin{equation}\label{eq:rank-criterion}
	P\mid F_m
	\quad\Longleftrightarrow\quad
	\rank(P)\mid m.
\end{equation}

If $b$ is odd, then
\begin{equation}\label{eq:square-free-root}
	F_b(x)=R_b(x)^2
\end{equation}
for a square-free polynomial $R_b$ \cite{HUNZIKER2004465}.
Consequently, if
\[
	H_b(x)=\gcd(R_b(x),R_b(x+1)),
\]
then
\begin{equation}\label{eq:base-nullity}
	d(b-1)=2\deg H_b.
\end{equation}
Translation by 1 exchanges the arguments defining $H_b$, so
$H_b(x+1)=H_b(x)$.
Since $H_b$ is square-free, translation pairs its roots as
$\{\alpha,\alpha+1\}$.
There are no fixed points, and hence $\deg H_b$ is even.
It follows that
\begin{equation}\label{eq:even-base-divisibility}
	4\mid d(b-1).
\end{equation}

We also use the root parameterization
\begin{equation}\label{eq:root-parameterization}
	F_m(\alpha)=0,\quad \alpha\neq0
	\quad\Longleftrightarrow\quad
	\alpha=\zeta+\zeta^{-1}
\end{equation}
for some non-trivial $m$th root of unity $\zeta$ in
$\overline{\F}_2$ \cite{HUNZIKER2004465}.

Finally, writing $n+1=2^s b$ with $b$ odd gives
\begin{equation}\label{eq:two-adic}
	d(n)
	=
	2^s d(b-1)
	+2(2^s-1)\mathbf{1}_{3\mid b}
\end{equation}
\cite{boylesResolutionSutner}.

\section{Bounding polynomial ranks}

\begin{lemma}\label{lem:rank-bound}
	Let $P\in\F_2[x]$ be a nonconstant square-free polynomial with
	$P(0)\neq0$.
	Then
	\[
		\rank(P)\leq2^{2\deg P}-1.
	\]
\end{lemma}

\begin{proof}
	Factor $P$ into distinct monic irreducibles,
	\[
		P=P_1\cdots P_t,
		\qquad
		e_i=\deg P_i.
	\]
	Choose a root $\alpha_i$ of $P_i$.
	Then $\alpha_i\in\F_{2^{e_i}}$ and $\alpha_i\neq0$.
	Let $\zeta_i$ be a root of
	\[
		z^2+\alpha_i z+1.
	\]
	This quadratic splits over an extension of degree at most two, so
	$\zeta_i\in\F_{2^{2e_i}}$.
	Furthermore, $\zeta_i\neq1$, since otherwise $\alpha_i=0$.
	If $m_i$ is the multiplicative order of $\zeta_i$, then
	\[
		m_i\mid 2^{2e_i}-1.
	\]
	By \eqref{eq:root-parameterization}, $\alpha_i$ is a root of
	$F_{m_i}$, and therefore
	\[
		\rank(P_i)\leq m_i\leq2^{2e_i}-1.
	\]

	Since the factors $P_i$ are distinct, the rank criterion gives
	\[
		\rank(P)
		=
		\mathop{\operatorname{lcm}}_{1\leq i\leq t}\rank(P_i).
	\]
	Hence
	\[
		\rank(P)
		\leq
		\prod_{i=1}^t(2^{2e_i}-1)
		<
		2^{2\sum_i e_i}
		=
		2^{2\deg P}.
	\]
	The rank is an integer, which gives the stated bound.
\end{proof}

\section{Compressing an occurrence}

\begin{theorem}[Odd-base compression]\label{thm:compression}
	Let $b$ be a positive odd integer and set $D=d(b-1)$.
	There is an odd divisor $b_0$ of $b$ such that
	\[
		d(b_0-1)=D,
		\qquad
		3\mid b_0\Longleftrightarrow3\mid b.
	\]
	If $D>0$, then $b_0$ may be chosen so that
	\[
		b_0\leq
		\begin{cases}
			2^D-1 & 3\nmid b,\\
			3(2^D-1) & 3\mid b.
		\end{cases}
	\]
	If $D=0$, one may take $b_0=1$ when $3\nmid b$ and $b_0=3$
	when $3\mid b$.
\end{theorem}

\begin{proof}
	The case $D=0$ follows from $d(0)=d(2)=0$, so suppose $D>0$.
	Write $F_b=R_b^2$ as in \eqref{eq:square-free-root}, and let
	\[
		H=\gcd(R_b(x),R_b(x+1)).
	\]
	The polynomial $H$ is square-free, and \eqref{eq:base-nullity} gives
	\[
		\deg H=\frac{D}{2}.
	\]
	Translation by 1 exchanges the two arguments of the GCD and has
	order two over $\F_2$, so
	\[
		H(x+1)=H(x).
	\]
	Also, $H(0)\neq0$ because $F_b(0)=1$ for odd $b$.

	Set $r=\rank(H)$.
	Since $H\mid F_b$, the rank criterion gives $r\mid b$.
	In particular, $r$ is odd.
	By definition, $H\mid F_r(x)$.
	Translating this divisibility and using $H(x+1)=H(x)$ gives
	\[
		H\mid F_r(x+1).
	\]
	Because $H$ is square-free and $r$ is odd, equation
	\eqref{eq:square-free-root} implies that $H$ divides both
	$R_r(x)$ and $R_r(x+1)$.
	Therefore
	\[
		d(r-1)\geq2\deg H=D.
	\]
	On the other hand, $r\mid b$, so Fibonacci divisibility gives
	\[
		d(r-1)\leq d(b-1)=D.
	\]
	Thus
	\begin{equation}\label{eq:rank-preserves-nullity}
		d(r-1)=D.
	\end{equation}

	Lemma \ref{lem:rank-bound} now gives
	\[
		r\leq2^{2\deg H}-1=2^D-1.
	\]
	If $3\nmid b$, take $b_0=r$.
	If $3\mid b$, take
	\[
		b_0=\operatorname{lcm}(r,3).
	\]
	In the second case, $r\mid b_0\mid b$.
	Combining \eqref{eq:rank-preserves-nullity} with Fibonacci
	divisibility gives
	\[
		D=d(r-1)\leq d(b_0-1)\leq d(b-1)=D.
	\]
	Hence $d(b_0-1)=D$.
	The construction preserves divisibility by 3, and the claimed bound
	follows from $b_0\leq3r$.
\end{proof}

\section{A finite witness theorem}

For a positive integer $N$, define the finite set
\[
	\mathcal{A}_N
	=
	\left\{
		(s,\epsilon,D):
		\begin{array}{l}
			s\geq0,\quad \epsilon\in\{0,1\},\quad
			D\geq0,\quad4\mid D,\\[2pt]
			N=2^sD+2\epsilon(2^s-1)
		\end{array}
	\right\}.
\]
For $(s,\epsilon,D)\in\mathcal{A}_N$, set
\[
	B_\epsilon(D)
	=
	\begin{cases}
		1 & D=0,\ \epsilon=0,\\
		3 & D=0,\ \epsilon=1,\\
		2^D-1 & D>0,\ \epsilon=0,\\
		3(2^D-1) & D>0,\ \epsilon=1.
	\end{cases}
\]
For $N>0$, the defining equation gives $2^s\leq N+2$.
Thus $\mathcal{A}_N$ is finite and can be enumerated effectively.

\begin{theorem}[Finite witness bound]\label{thm:finite-witness}
	Let $N>0$.
	If $d(n)=N$ for some $n$, then $\mathcal{A}_N$ is nonempty and
	there is another occurrence $d(n_0)=N$ with
	\begin{equation}\label{eq:refined-bound}
		n_0
		\leq
		\max_{(s,\epsilon,D)\in\mathcal{A}_N}
		\left(2^sB_\epsilon(D)-1\right).
	\end{equation}
	In particular,
	\begin{equation}\label{eq:coarse-bound}
		n_0\leq3(N+2)2^N-1.
	\end{equation}
\end{theorem}

\begin{proof}
	Suppose $d(n)=N$, and write $n+1=2^s b$ with $b$ odd.
	Set
	\[
		D=d(b-1),
		\qquad
		\epsilon=\mathbf{1}_{3\mid b}.
	\]
	Equation \eqref{eq:two-adic} gives
	\[
		N=2^sD+2\epsilon(2^s-1).
	\]
	Moreover, \eqref{eq:even-base-divisibility} gives $4\mid D$, so
	$(s,\epsilon,D)\in\mathcal{A}_N$.

	Apply Theorem \ref{thm:compression} to obtain $b_0$.
	The divisibility condition on 3 is preserved, so another application
	of \eqref{eq:two-adic} gives
	\[
		d(2^s b_0-1)=N.
	\]
	The bound on $b_0$ proves \eqref{eq:refined-bound}.

	It remains to prove the coarse bound.
	For $\epsilon=0$, the equality $N=2^sD$ and $N>0$ give
	$2^s\leq N$.
	For $\epsilon=1$,
	\[
		N+2=2^s(D+2),
	\]
	so $2^s\leq N+2$.
	In both cases $D\leq N$ and
	\[
		B_\epsilon(D)\leq3\cdot2^D.
	\]
	Therefore
	\[
		2^sB_\epsilon(D)-1
		\leq3(N+2)2^N-1,
	\]
	as desired.
\end{proof}

\begin{corollary}\label{cor:decidable}
	Whether an integer occurs as $d(n)$ is decidable by a finite
	computation.
\end{corollary}

\begin{proof}
	Negative integers are impossible because nullities are nonnegative,
	the value 0 occurs, and every positive odd value is impossible because
	$d(n)$ is always even.
	For a positive even $N$, compute the coarse bound
	\eqref{eq:coarse-bound} and evaluate the finitely many values through
	that bound using \eqref{eq:nullity}.
\end{proof}

\section{Computational complexity}

The complexity depends strongly on how the input is encoded.
The $O(n\log^2 n)$ evaluation bound of
\cite{GoldwasserKlostermeyerTrapp1997} is polynomial in the numerical
board dimension $n$, equivalently in a unary encoding of $n$.
For binary-encoded $n$, the same algorithm is exponential in the input
length.

\begin{corollary}\label{cor:complexity}
	For $N>0$, using the bound in Theorem \ref{thm:finite-witness},
	$\operatorname{IMAGE}_d$ can be decided in time $2^{O(N)}$, up to
	polynomial factors in the target nullity $N$.
	If $N$ is encoded in binary with length
	$L=\lceil\log_2(N+1)\rceil$, this gives the
	double-exponential upper bound
	\[
		2^{2^{O(L)}}.
	\]
\end{corollary}

\begin{proof}
	Set
	\[
		B(N)=3(N+2)2^N-1.
	\]
	By Theorem \ref{thm:finite-witness}, it suffices to evaluate $d(n)$
	for $0\leq n\leq B(N)$.
	Applying the published evaluator independently at every index takes
	\[
		\sum_{n\leq B(N)}O(n\log^2 n)
		=
		O(B(N)^2\log^2 B(N))
		=
		2^{O(N)}.
	\]
	Since $N<2^L$, this becomes $2^{2^{O(L)}}$ in terms of the binary
	input length.
\end{proof}

\begin{remark}
	This is only a coarse upper bound.
	It supplies neither a matching lower bound nor an NP verification
	procedure, and it does not rule out substantially smaller witnesses
	or more succinct certificates.
	We are unaware of a published hardness result for
	$\operatorname{IMAGE}_d$, or for evaluating $d(n)$ when $n$ is
	encoded in binary.
\end{remark}

\section{The impossible nullity 34}

\begin{corollary}\label{cor:34-bound}
	If $d(n)=34$ for some $n$, then $d(n_0)=34$ for some
	\[
		n_0\leq393197.
	\]
\end{corollary}

\begin{proof}
	The equation defining $\mathcal{A}_{34}$ is
	\[
		34=2^sD+2\epsilon(2^s-1),
		\qquad
		4\mid D.
	\]
	If $\epsilon=0$, the only integral possibilities are
	$(s,D)=(0,34)$ and $(1,17)$, neither of which has $4\mid D$.
	If $\epsilon=1$, then
	\[
		36=2^s(D+2).
	\]
	The integral possibilities give $D=34,16,7$ for $s=0,1,2$,
	respectively.
	Hence the only admissible triple is
	\[
		(s,\epsilon,D)=(1,1,16).
	\]

	In the proof of Theorem \ref{thm:compression}, the rank $r$ is odd
	and satisfies
	\[
		r\leq2^{16}-1=65535.
	\]
	If $3\mid r$, then $b_0=r\leq65535$.
	If $3\nmid r$, then oddness gives $r\leq65533$, and therefore
	\[
		b_0=3r\leq196599.
	\]
	Thus in all cases
	\[
		n_0=2b_0-1\leq393197.
	\]
\end{proof}

\begin{theorem}\label{thm:34-impossible}
	For every $n\geq0$,
	\[
		d(n)\neq34.
	\]
\end{theorem}

\begin{proof}
	Corollary \ref{cor:34-bound} does not require computing every
	nullity through 393197.
	The unique admissible triple shows that 34 occurs if and only if
	there is an odd multiple $b$ of 3 such that
	\[
		d(b-1)=16.
	\]
	The compression theorem shows that it suffices to check these $b$
	through 196599.
	There are 32,767 such candidates.
	The data in \path{code/serialization/b159257.txt} through $n=100000$
	checks the first 16,667 of them and finds none.
	We computed the remaining 16,100 values exactly using packed
	Fibonacci polynomials and the \texttt{GF2X} implementation in
	NTL 11.6.0 \cite{ShoupNTL}.
	The conversion and exact GCD calculation are implemented in
	\path{code/ai_prototypes/wieferich_check.py} and
	\path{code/ai_prototypes/ntl_gf2x_gcd.cpp}.
	The exhaustive driver is
	\path{code/ai_prototypes/verify_nullity_34.py}; it may be run with
	\begin{center}
		\texttt{cd code \&\& python -m ai\_prototypes.verify\_nullity\_34}.
	\end{center}
	The canonical SHA-256 digest of the 16,100 newly computed pairs is
	\path{64cc5c41bf1bf197c833f695e3aaf3308b4154b358667f1a4f79b40a0a2458ab}.
	None has $d(b-1)=16$.
	A separate computation with the Python polynomial-GCD implementation
	agreed on 104 distributed sample values.
	Thus no candidate remains, and 34 is impossible.
\end{proof}

\section{Discussion}

Theorem \ref{thm:finite-witness} gives an a priori finite test for every
candidate nullity.
Its coarse bound is not intended to be computationally sharp.
The useful part of the argument is that an occurrence at an arbitrarily
large odd base can be replaced by one whose base is controlled by the
degree of its common factor.

For a particular target $N$, one should first enumerate the admissible
triples in $\mathcal{A}_N$ and then apply the corresponding bounds from
Theorem \ref{thm:compression}.
Further improvements are possible by classifying the irreducible factors
of $H$ and their ranks, as in the proof that 26 is impossible, instead of
using the general estimate in Lemma \ref{lem:rank-bound}.

The result for 34 demonstrates the role of the witness theorem.
Without Theorem \ref{thm:finite-witness}, absence through any finite
cutoff would provide only evidence.
With the theorem, exhausting the bounded candidate set proves global
impossibility.

\bibliographystyle{amsalpha}
\bibliography{refs.bib}

\end{document}
